\appendixsection{附录 关于线性表示的补充}

下面的命题推广了 Chap. 1, §3, no. 8 的 Prop. 13。

命题 1. 设 K 是一个交换域，A 是一个 K-代数，V 与 W 是两个作为 K 上有限维向量空间的左 A-模。若存在 K 的一个扩张 L，使得 \((A \otimes_{K} L)\)-模 \(V \otimes_{K} L\) 与 \(W \otimes_{K} L\) 同构，则 A-模 V 与 W 同构。

\begin{enumerate}
\def\labelenumi{\alph{enumi})}
\item
  先设 L 是 K 的次数为 n~的有限扩张。由于 \(V \otimes_{K} L\) 与 \(W \otimes_{K} L\) 作为 \((A \otimes_{K} L)\)-模同构，它们作为 A-模也同构；而作为 A-模，它们分别同构于 \(V^{n}\) 与 \(W^{n}\)。V 与 W 是有限长度的 A-模；于是 V（相应地，W）是子模族 \((\mathrm{M}_{i}^{r_{i}})_{1 \leqslant i \leqslant p}\)（相应地，\((\mathrm{N}_{j}^{s_{j}})_{1 \leqslant j \leqslant l}\)）的直和，使得诸 \(M_{i}\)（相应地，\(N_{j}\)）不可分解，且指标不同的两个 \(M_{i}\)（相应地，\(N_{j}\)）互不同构（Algebra, Chap. VIII, §2, no. 2, Th. 1）。于是 \(V^{n}\)（相应地，\(W^{n}\)）是诸 \(M_{i}^{nr_{i}}\)（相应地，\(N_{j}^{ns_{j}}\)）的直和；由此（loc. cit.）可知 p = q，且在适当置换诸 \(N_{j}\) 之后，\(M_{i}\) 同构于 \(N_{i}\) 且 \(nr_{i}\) 等于 \(ns_{i}\) 对 \(1 \leqslant i \leqslant p\) 成立。于是 V 同构于 W。
\item
  设 K 是一个无限域。假设蕴含 V 与 W 在 K 上有相同的维数。设 \((e_{i})_{1\leqslant i\leqslant m}\) 与 \((e_{i}^{\prime})_{1\leqslant i\leqslant m}\) 是 V 与 W 在 K 上的基，\((a_{\lambda})\) 是 A 在 K 上的基。一个同构 \(u:V\otimes_{K}L\to W\otimes_{K}L\) 是一个双射的 L-线性映射，同时也是一个 \((A\otimes_{K}L)\)-同态，换句话说，它满足条件
\end{enumerate}

\[
a _ {\lambda} u \left(e _ {i}\right) = u \left(a _ {\lambda} e _ {i}\right) \quad \text { for   all } \lambda \text { and   all } i.\tag{1}
\]

置 \(a_{\lambda}e_{i} = \sum_{j}\gamma_{\lambda ij}e_{j}, a_{\lambda}e_{i}' = \sum_{j}\gamma_{\lambda ij}'e_{j}'\)，其中诸 \(\gamma_{\lambda ij}\) 与 \(\gamma_{\lambda ij}'\) 属于 K；并置 \(u(e_i) = \sum_{j}\xi_{ij}e_j'\)，其中诸 \(\xi_{ij}\) 属于 L。条件 (1) 可写成

\[
\sum_ {j} \xi_ {i j} \gamma_ {\lambda j k} ^ {\prime} = \sum_ {j} \gamma_ {\lambda i j} \xi_ {j k}\tag{2}
\]

对所有 \(\lambda, i, k\) 成立。按假设，齐次线性方程组 (2) 有一个解 \((\xi_{ij}) \in L^{m^2}\) 满足 \(\det(\xi_{ij}) \neq 0\)。由于方程组 (2) 的系数属于 K，我们知道（Algebra, Chap. II, §8, no. 5, Prop. 6）该方程组在 \(K^{m^2}\) 中也承认非平凡解；设 E 是 \(\mathbf{M}_{m}(\mathrm{K})=\mathrm{K}^{m^{2}}\) 中由这些解构成的、不退化为零的向量子空间。设 \((c_{l})_{1\leqslant l\leqslant p}\) 是 E 的基，置 \((\xi_{ij})=\sum_{l}\eta_{l}c_{l}\)（对任何矩阵 \((\xi_{ij})\in\mathrm{E}\)）；于是 \(\det(\xi_{ij})\) 是系数在 K 中的多项式 \(\mathrm{P}(\eta_{1},\ldots,\eta_{p})\)。另一方面，我们知道（loc. cit.）(2) 在 \(L^{m^{2}}\) 中的解形如 \(\sum_{l}\zeta_{l}c_{l}\)，这里 \(\zeta_{l}\in L\)；对这样的解，\(\det(\xi_{ij})\) 等于 \(\mathrm{P}(\zeta_{1},\ldots,\zeta_{p})\)。据此，若 \(\mathrm{P}(\eta_{1},\ldots,\eta_{p})=0\) 对所有 \(\eta_{1},\ldots,\eta_{p}\in K\) 都成立，则由于 K 无限，P 的系数全为零；于是 \(\mathrm{P}(\zeta_{1},\ldots,\zeta_{p})=0\) 将对所有 \(\zeta_{1},\ldots,\zeta_{p}\in L\) 成立，这与我们的假设相反。因此可以找到一个矩阵 \((\xi_{ij})\in\mathrm{E}\) 使得 \(\det(\xi_{ij})\neq0\)，而相应的线性映射 \(V\to W\) 是一个同构。

\begin{enumerate}
\def\labelenumi{\alph{enumi})}
\setcounter{enumi}{2}
\tightlist
\item
  一般情形。设 \(\Omega\) 是 \(L\) 的一个代数闭扩张，\(K_0\) 是 \(K\) 在 \(\Omega\) 中的代数闭包。假设蕴含 \(V \otimes_K \Omega\) 与 \(W \otimes_K \Omega\) 作为（\(A \otimes_K \Omega\)）-模同构。由于 \(K_0\) 无限，\(b\) ) 表明 \(V \otimes_K K_0\) 与 \(W \otimes_K K_0\) 作为（\(A \otimes_K K_0\)）-模同构。沿用 \(b\) ) 中的记号，方程组 (2) 有一个解（\(\xi_{ij}\)）\(\in K_0^{m^2}\) 满足 \(\det(\xi_{ij}) \neq 0\)。但诸 \(\xi_{ij}\) 都属于某个代数扩张 \(K_1\)，它在 \(K\) 上的次数有限。（\(A \otimes_K K_1\)）-模 \(V \otimes_K K_1\) 与 \(W \otimes_K K_1\) 同构，于是利用 \(a\) ) 即完成证明。
\end{enumerate}

命题 2（Maschke）. 设 A 是一个有单位元的环，E 是一个左 A-模，F 是 E 的一个直和项，G 是一个阶为 q 的有限群，\(\rho\) 是 G 在 E 上的一个线性表示。假设 q.1 在 A 中可逆且 F 在 G 下稳定。则存在 F 在 E 中的一个在 G 下稳定的补。

设 p 是 E 到 F 上的一个投射。对所有 \(x \in E\)，置

\[
f (x) = q ^ {- 1} \sum_ {s \in G} \rho (s) ^ {- 1} p (\rho (s) x).
\]

我们有 \(f(x) \in \mathbf{F}\)，且 \(f(y) = y\) 对所有 \(y \in \mathbf{F}\) 都成立，故 \(f\) 是 \(\mathbf{E}\) 到 \(\mathbf{F}\) 上的一个投射。另一方面，若 \(t \in \mathbf{G}\)，则

\[
\begin{array}{r l} \rho (t) f (x) & = q ^ {- 1} \sum_ {s \in G} \rho (s t ^ {- 1}) ^ {- 1} p (\rho (s) x) \\ & = q ^ {- 1} \sum_ {s \in G} \rho (s) ^ {- 1} p (\rho (s t) x) \\ & = f (\rho (t) x). \end{array}
\]

于是 \(f\) 与 \(\rho(\mathrm{G})\) 交换，故 \(\operatorname{Ker}(f)\) 是 \(\mathrm{F}\) 的一个在 \(\mathrm{G}\) 下稳定的补。

系。设 G 是一个阶为 q 的有限群，K 是一个特征不整除 q 的交换域。则 G 在 K 上的群代数是半单的。

事实上，由 Prop. 2，这个代数上的每个模都是半单的。

命题 3. 设 A 是一个交换环，M 是一个 A-模，G 是一个作用在 M 上的有限群，A' 是一个 A-模。假设 G 的阶 q 在 A 中可逆。设 \(M^{G}\) 是 M 中在 G 下不变的元素的集合。则从 \(M^{G} \otimes_{A} A'\) 到 \(M \otimes_{A} A'\) 的典范同态定义了从 \(M^{G} \otimes_{A} A'\) 到模 \((M \otimes_{A} A')^{G}\)（即 \(M \otimes_{A} A'\) 中在 G 下不变的元素所成的模）上的一个同构。

事实上，设 \(\mathbf{Q}\) 是 \(\mathbf{M}\) 到 \(\mathbf{M}^{\mathbf{G}}\) 上的、由 \(\mathbf{Q}(x) = q^{-1}\sum_{g\in \mathbf{G}}g(x)\)（对所有 \(x\in \mathbf{M}\)）所定义的投射。若 \(i\) 表示 \(\mathbf{M}^{\mathbf{G}}\) 到 \(\mathbf{M}\) 中的典范单射，则 \(\mathbf{Q}\circ i\) 是 \(\mathbf{M}^{\mathbf{G}}\) 的恒同映射，故 \((\mathbf{Q}\otimes 1_{\mathbf{A}'})\circ (i\otimes 1_{\mathbf{A}'})\) 是 \(\mathbf{M}^{\mathbf{G}}\otimes_{\mathbf{A}}\mathbf{A}'\) 的恒同映射。由于 \(\mathbf{Q}\otimes 1_{\mathbf{A}'} = q^{-1}\sum_{g\in \mathbf{G}}(g\otimes 1_{\mathbf{A}'})\)，\(i\otimes 1_{\mathbf{A}'}\) 的像为 \((\mathbf{M}\otimes \mathbf{A}')^{\mathbf{G}}\)。另一方面，由前面所说的内容，\(i\otimes 1_{\mathbf{A}'}\) 是单射。

注。前一命题特别适用于 \(A'\) 是一个 \(A\)-代数的情形。此时，\(M^G \otimes_A A'\) 是一个 \(A'\)-子模，含于 \(M \otimes_A A'\)。

